Para confirmar que a tensão da eq. (3.27) satisfaz a condição de consistência da eq. (3.24), realiza-se a verificação dimensional explícita.

\textbf{Dimensões relevantes para estruturas:}
\begin{itemize}
\item Tensão $\sigma$: $[\sigma] = M\,L^{-1}\,T^{-2}$ (força por área)
\item Momento fletor $M_b$: $[M_b] = M\,L^{2}\,T^{-2}$
\item Distância ao centroide $y$: $[y] = L$
\item Segundo momento de área $I$: $[I] = L^{4}$
\item Módulo de elasticidade $E$: $[E] = M\,L^{-1}\,T^{-2}$
\end{itemize}

\textbf{Substituição da tensão de flexão $\sigma = M_b\,y/I$:}
\[
[\sigma] = \frac{[M_b]\,[y]}{[I]} = \frac{M\,L^{2}\,T^{-2}\cdot L}{L^{4}} = \frac{M\,L^{3}\,T^{-2}}{L^{4}} = M\,L^{-1}\,T^{-2}.\quad\checkmark
\]

O resultado $M\,L^{-1}\,T^{-2}$ é exatamente a dimensão de tensão exigida pela eq. (3.24). Portanto, a tensão da eq. (3.27) satisfaz a condição de consistência dimensional.

